\begin{tablalafrox}{Anexo 7.F — Supuestos de planificación y dependencias para la gestión de riesgos. Fuente: elaboración propia.}{tab_7F1}
  {Y Y Y Y Y}
  {\cab{ID} & \cab{Riesgo / supuesto} & \cab{EDT / hitos} & \cab{Responsable de tratamiento} & \cab{Disparador o control}}
  P7-01 & Validación de productividad y tamaños HH por clase & Los 222 paquetes; T-15, sección 4 & JP y líderes de frente & Se valida con el equipo en la línea base (H1), con trazabilidad a T-12, cantidades y ensayos. Se recalcula si la demanda supera la capacidad \\
  P7-02 & Cronograma por actividad depende de tamaños y equipos supuestos & D-01–D-34; H2/\allowbreak{}H3/\allowbreak{}H4/\allowbreak{}H5/\allowbreak{}H9/\allowbreak{}H10 & JP/ARQ & T-15, secciones 5 y 6, programa 564 actividades con dependencias, revisión Art. 18.3 y nivelación; reservas de 6 a 35 días hábiles por hito. Se valida con el equipo en la línea base (H1) y se repite el cálculo. Se escala y replanifica cuando la desviación proyectada supera la mitad de la reserva \\
  P7-03 & Solapamientos compiten por especialistas & 4.2.1/\allowbreak{}4.2.2, 3.5, 4.3.1; meses 13–15/\allowbreak{}19–20 & Líder DES y CAL & Mantener 256 HH DES + 128 HH CAL/mes protegidas E1, sin préstamo a F4; asignar personas nominales \\
  P7-04 & Fecha efectiva cambia congelamientos & 1.1.3, 4.1; V-12/\allowbreak{}H6/\allowbreak{}H11/\allowbreak{}H7/\allowbreak{}H12 & JP/CLIENTE & Confirmar fecha y transformar meses relativos en calendario; no iniciar corte en fechas prohibidas \\
  P7-05 & Ola/cadena no lista antes de cuatro semanas de cierre & 4.2.1/\allowbreak{}4.3.1, 7.3, 3.6.5/\allowbreak{}6 & IMP/Comercial/Operaciones & Registro de todo el alcance, activación antes del tramo final, cero incidentes críticos/altos; no reemplazar alcance por muestra \\
  P7-06 & Reversión manual no soporta despacho o falla DTE & 4.1.2/\allowbreak{}3.3.2; H6/\allowbreak{}H11 & Operaciones/ARQ/SRE & Ensayo de 96 despachos, objetivo total 40 minutos, final antes de 05:30; contingencia ERP aprobada \\
  P7-07 & Retiro temprano de acompañamiento o capacidad de mesa insuficiente & 4.2.2/\allowbreak{}4.3.2/\allowbreak{}8.1.1/\allowbreak{}8.1.2/\allowbreak{}8.1.5 & IMP/SRE & Decremento sólo con acta e indicadores sostenidos; mesa con las posiciones del SD4 y SOC 24\ensuremath{\times}7 en T-15; medir abandono y resolución al primer contacto y calibrar Erlang A (T-15 §5.6) \\
  P7-08 & Compra/sala/borde incumple H3 & 5.1.2/\allowbreak{}6.1/\allowbreak{}6.3/\allowbreak{}6.6.3 & SRE & LafroX emite la orden de compra en el mes 2. La sala se instala en el mes 3 y se recibe en el mes 4, antes de racks y configuración. El borde de los CD entra en servicio en el mes 5 para el H3 \\
  P7-09 & Reserva de stock/retención duplica o confirma sin acuse & 3.3.6/\allowbreak{}3.4.2/\allowbreak{}3.4.6/\allowbreak{}3.8.1/\allowbreak{}3.8.3 & ARQ/DES/CAL & La coordinación de reserva y retención (SD4, apartado 4.1.4.4) se valida antes del H4 con AL-STOCK-01/\allowbreak{}AL-ACT-01: concurrencia, corte, reintento, acuse durable y autoridad por época \\
  P7-10 & Tráfico adicional o recuperación excede diseño & 3.8.4/\allowbreak{}3.9.3; SD4 física/DR & ARQ/SRE & AL-STOCK-01 mide la proporción de líneas con más de una retención y 3.8.4/\allowbreak{}3.9.3 verifican carga y drenaje conforme al Anexo 4-I, Tabla A.10, y Anexo 4-W, Tablas A.32 y A.33, sin sumar la coordinación al drenaje. El riesgo residual de falla de los tres caminos seguida de destrucción del sitio (SD4 4.3.2.4; R8-05) se verifica con AL-DR-01 y sus alarmas de 5 y 15 min, reposición del enlace, preemisión de guías y NAS WORM \\
  P7-11 & Control de conservación del precio pactado & 1.2.1/\allowbreak{}3.4.6; SD2 S-09, RF-03.11/\allowbreak{}12 y SD3 RNG-08 & Comercial/JP & Probar cambio de lista entre pedido y despacho y detectar diferencias ERP \\
\end{tablalafrox}
